\section{From curvature rarity to full-measure box windows}
\label{fbb:section}

The two preceding probability estimates can be joined without
assuming that a boundary field is typically close to the identity
in an unfixed gauge. The resulting windows are defined by an
explicit local tree gauge and a conditional Wilson minimizer.
They still differ from the nonlinear coarse-average windows in
the homogeneous renormalization construction.

\begin{lemma}[A boundary-only gauge at fixed geometry]
\label{fbb:boundarygauge}
For the fixed four-dimensional cubical box \(\Omega\), choose a
spanning tree \(\mathcal T_\partial\) of its boundary edge graph
and a root vertex. There is a finite constant \(L_R\) such that
every boundary field has a boundary gauge, determined only by its
boundary edges, with
\[
 V_e=I\quad(e\in\mathcal T_\partial),\qquad
 \max_{e\in E_\partial}\norm{V_e-I}_\tau
 \leq L_R\max_{p\subset\partial\Omega}\norm{U_p-I}_\tau.
\]
The estimate concerns boundary plaquettes, not plaquettes through
the interior. No assertion uniform in \(R\) is made.
\end{lemma}
\begin{proof}
Fix the root gauge to \(I\) and gauge the boundary tree edges to \(I\)
by recursion. Every remaining boundary edge becomes the holonomy
of its fixed tree fundamental loop. The boundary cubical complex is
homeomorphic to \(S^3\); its two-skeleton has trivial fundamental
group. Each of the finitely many fundamental loops is therefore a
finite product of conjugates of elementary boundary plaquette
loops and their inverses, as a word in the graph edge generators.
This assertion follows directly from the cellular presentation of
the fundamental group: the plaquette boundary words normally
generate all relations in the graph fundamental group. Choose one
such finite word once for each non-tree edge and let \(L_R\) be the
largest number of plaquette factors used.

For unitary matrices \(A_j\), telescoping gives
\(\norm{\prod_j A_j-I}_\tau\leq\sum_j\norm{A_j-I}_\tau\).
Conjugation and inversion preserve this distance. Applying this
inequality to the chosen word proves the bound, regardless of
the size of its conjugating holonomies.
\end{proof}

Choose \(\varepsilon_R>0\) so small that
\(\sqrt2 L_R\varepsilon_R\) lies inside the small boundary
neighborhood in Theorem~\ref{fbw:theorem}. We may and do take that
neighborhood to be a conjugation-invariant product ball. Call a
box boundary good when
\[
 h(U_p)<\varepsilon_R^2\quad\hbox{for every }p\subset\partial\Omega.
\]
On a good boundary, apply the boundary-only gauge in
Lemma~\ref{fbb:boundarygauge}; extend its vertex values into the
box arbitrarily, then perform the rooted forest gauge of
Section~\ref{fbw:section}, keeping the transformed boundary fixed.
The final representative is independent of that arbitrary extension:
two extensions differ by an interior gauge, which the unique forest
normalization removes. It is also the representative obtained by
gauging the full spanning tree
\(\mathcal T_\partial\cup\mathcal F\) to \(I\).

Define the box bad event \(D_\Omega(g,T)\) as follows:
the boundary is bad, or the boundary is good and the resulting
interior representative lies outside \(W_{V,g,T}\). All these maps
are measurable. The event is gauge invariant. Indeed, an arbitrary
original gauge transformation changes the full tree representative
only by constant conjugation at its root. The conditional minimizer
transforms by that same conjugation, by uniqueness, and the local
logarithm and edge-max norm respect conjugation.

\begin{theorem}[Prescribed local box defects in the full Wilson measure]
\label{fbb:boxes}
Fix \(R,b\), with \(b\geq3\), and consider a periodic four-dimensional
lattice with side lengths divisible by \(2b\). Let
\(\Omega_1,\ldots,\Omega_m\) be embedded translates of the fixed box
with pairwise disjoint vertex sets. Put
\[
 Q_b=b^4\binom42,\qquad
 N_\partial=\#\{p:p\subset\partial\Omega\}=24R^3,
\]
and let \(z_b(g,\varepsilon_R)\) be the bad-cell factor of
Section~\ref{fwp-section}. There are constants \(C,c,r,g_0>0\),
independent of the total torus size and of the chosen boxes, such
that
\begin{equation}
 \mu_g\left(\bigcap_{i=1}^m D_{\Omega_i}(g,T)\right)
 \leq \left[
       C e^{-cT^2}
       +N_\partial z_b(g,\varepsilon_R)^{1/Q_b}
       \right]^m
 \quad(0<g\leq g_0,\ 1\leq T\leq r/g).
 \label{fbb:joint}
\end{equation}
In particular, for every \(a>1/2\), the bracket with
\(T(g)=(\log(1/g))^a\) is smaller than every fixed power of \(g\)
as \(g\downarrow0\).
\end{theorem}
\begin{proof}
Condition on every edge outside the union of the boxes' interior
edge sets \(E_{\circ,i}\). This includes all box boundary edges.
Every plaquette meeting \(E_{\circ,i}\) lies wholly in \(\Omega_i\).
Disjointness of the vertex sets therefore ensures that no plaquette
meets two integrated edge sets. The conditional law factorizes
exactly into the independent conditional Wilson laws of the
individual boxes; all remaining action terms are constant under
these integrations.

Write \(B_i\) for a bad boundary and \(F_i\) for the interior-window
failure on a good boundary. Each \(B_i\) is measurable under the
conditioning. The boundary gauge is also fixed under it, and it
preserves the conditional product Haar law and the action.
Theorem~\ref{fbw:theorem} consequently gives
\[
 \mathbb P(F_i\mid\hbox{conditioned edges})
 \leq a_g\,\mathbf1_{B_i^c},\qquad a_g=Ce^{-cT^2}.
\]
The \(F_i\)'s are conditionally independent.

For any subset \(J\) of boxes, a union bound over one offending
boundary plaquette in each box, followed by the prescribed-plaquette
bound of Section~\ref{fwp-section}, gives
\begin{equation}
 \mu_g\left(\bigcap_{i\in J}B_i\right)
 \leq \left[N_\partial
              z_b(g,\varepsilon_R)^{1/Q_b}\right]^{|J|}.
 \label{fbb:boundaryrare}
\end{equation}
The selected plaquettes are distinct because the boxes have disjoint
vertex sets. This is a correlated full-measure estimate; independence
of the boundaries is not used.

Expand \(\bigcap_i(B_i\cup F_i)\) over \(J\), the boxes assigned
the boundary-bad alternative. Conditional factorization bounds
the contribution for \(J\) by
\(a_g^{m-|J|}\mu_g(\bigcap_{i\in J}B_i)\), after dropping other
good-boundary indicators. Substituting
\eqref{fbb:boundaryrare} and summing the binomial expression proves
\eqref{fbb:joint}.

Finally, the explicit bad-cell exponent has fixed threshold
\(\varepsilon_R\) and fixed \(b\):
\[
 z_b(g,\varepsilon_R)^{1/Q_b}
 \leq \exp\left[-\frac{\varepsilon_R^2}{Q_b g^2}
                       +C_{b,G}\log(1/g)+C_{b,G}\right]
\]
for sufficiently small \(g\). This and
Corollary~\ref{fbw:logwindow} prove the last assertion.
\end{proof}

\paragraph{What this joins, and what it does not.}
The theorem gives joint rarity of actual, gauge-invariant local
coordinate defects under the complete bare Wilson measure. It
combines a correlated boundary estimate with conditional interior
normalization and does not replace the Wilson measure by a Gaussian
law. For a fixed separated grid of such boxes, the usual connected-set
counting yields Peierls tails with the displayed bracket.

The representative and minimizer are those of a fixed conditional
box with all its boundary edges held fixed. They are not identified
with the averaged, constrained minimizer and independent retained
cutoffs of the existing multiscale map. No derivative regularity,
complex-source estimate, growing-box constant, or post-renormalization
reflection positivity follows from this theorem. Thus it is a
concrete full-measure local bridge, while the multiscale mixed-history
induction remains an additional theorem to prove.
